\documentclass[11pt]{article}

\usepackage[T1]{fontenc}
\usepackage[utf8]{inputenc}
\usepackage[english]{babel}
\usepackage[top=2.6cm, bottom=2.5cm, left=2.4cm, right=2.4cm]{geometry}
\usepackage{amsmath, amssymb, amsfonts, mathtools, bm, nicefrac}
\usepackage{booktabs}
\usepackage{array}
\usepackage{longtable}
\usepackage{microtype}
\usepackage{xcolor}
\usepackage{fancyhdr}
\usepackage[colorlinks=true, linkcolor=blue, citecolor=blue, urlcolor=blue]{hyperref}

% Matches commands.sty in Tensor_hyper_contraction__THC_.
\newcommand{\mat}[1]{{\mathbf{#1}}}
\newcommand{\parth}[1]{\left(#1\right)}
\newcommand{\sqrbrckts}[1]{\left[#1\right]}

\newcommand{\Dcal}{\mathcal{A}}
\newcommand{\Qcal}{\mathcal{Q}}
\newcommand{\Ncal}{\mathcal{N}}
\newcommand{\eb}{\tilde{\epsilon}}
\newcommand{\Had}{\circ}
\newcommand{\code}[1]{\texttt{\small #1}}

\setlength\parindent{0pt}
\setlength\parskip{8pt}

\pagestyle{fancy}
\setlength{\headheight}{14pt}
\lhead{\normalsize\textrm{Marcus Allen}}
\rhead{THC dRPA: attempted and replaced}
\renewcommand{\headrulewidth}{0.4pt}

\title{\vspace{-1.5cm}Tensor hypercontraction for the dRPA moments:\\
what was attempted, and what replaces it}
\author{Marcus Allen}
\date{26 August 2026}

\begin{document}
\maketitle
\vspace{-1.0cm}

\begin{abstract}
\noindent
Between September 2025 and March 2026 a tensor-hypercontracted dRPA was written eleven
times. Every version shares one architecture and one difficulty: the particle--hole
denominator is not separable in the occupied and virtual indices, so the pair space was
partitioned by the size of the denominator relative to the quadrature point and each region
given a different expansion. This note sets out those expansions explicitly, shows why the
partition was forced --- a separable builder can only address a rectangle in the pair space,
while the region it needs is a staircase --- and records where the accuracy was lost. It
then gives the replacement. The partition is an artefact of an oscillatory inner quadrature
rather than of the factorisation; representing the denominator on a fitted grid of
exponentials shared by every pole removes the quadrature, and with it the partition, the
threshold parameters and about nine hundred lines. A companion note derives the replacement
in full; here the emphasis is on the correspondence between the two.
\end{abstract}

\section{The object, in one paragraph}

Let $D_{ia} = \epsilon_a - \epsilon_i > 0$ be the particle--hole gaps and $\mat{V}$ the
vertex. The even density--density moments follow from
\begin{align}
    \Ncal^{(m)}
    = \frac{2}{\pi} \int_0^\infty
      \parth{\mat{I} + 4\Qcal^{(1)}(z)}^{-1} \Qcal^{(m+1)}(z) \,\mathrm{d}z ,
    \qquad
    \Qcal^{(k)}(z) = \mat{V}\, \frac{\mat{D}^{k}}{\mat{D}^2 + z^2}\, \mat{V}^T ,
    \label{eq:compact}
\end{align}
the odd ones from the shift-independent $\Dcal^{(m)} = \mat{V}\mat{D}^m\mat{V}^T$, and both
families feed one auxiliary-space recursion. Under tensor hypercontraction the vertex
factorises across the pair, $V_{P,ia} \to X_{iP} X_{aP}$ (times a Cholesky factor of the
Coulomb matrix), so
\begin{align}
    \Qcal^{(k)}_{PQ} = \sum_{ia} \parth{X_{iP} X_{iQ}} \parth{X_{aP} X_{aQ}}\, f(D_{ia}) ,
    \qquad f(D) = \frac{D^k}{D^2+z^2} .
    \label{eq:contract}
\end{align}
The two bracketed factors separate; $f(D_{ia})$ does not. Everything below is about that
one obstruction.

\section{The architecture every previous version shared}

All eleven versions do the same thing. An outer Clenshaw--Curtis or Gauss--Laguerre
quadrature on $[0,\infty)$ discretises \eqref{eq:compact}, with its grid spacing optimised
per calculation against an exactly evaluable diagonal integral; at each point $z$ the
auxiliary-space objects $\Qcal^{(m+1)}(z)$ are assembled, and the accumulation
\begin{align}
    \Ncal^{(m)} \;\mathrel{+}= \; w_z \,\frac{2}{\pi}
    \parth{\mat{I} + 4\Qcal^{(1)}(z)}^{-1} \Qcal^{(m+1)}(z)
    \label{eq:accum}
\end{align}
is performed. Equation \eqref{eq:accum} is correct and is kept: the implementation
without hypercontraction, \code{rpa\_new\_clean}, does exactly this and works. What varies between
versions is only how $\Qcal^{(m+1)}(z)$ is obtained from \eqref{eq:contract}, and that is
where the complexity lives.

Because $f(D)$ is not separable, the pair space was partitioned by the ratio
\begin{align}
    r_{ia} = \frac{D_{ia}}{z}
\end{align}
into three regions, each with its own expansion. The three are set out next; each is stated
first as a scalar identity in $D$ and then as the auxiliary-space object the code builds.

\subsection{\texorpdfstring{The series branch: $r \ll 1$}{The series branch: r much less than 1}}

For $D < z$ the resolvent is expanded as a geometric series in $r^2$,
\begin{align}
    \frac{D^{2\rho+1}}{D^2+z^2}
    \;=\; \sum_{c \ge 0} (-1)^c \, \frac{D^{2\rho+1+2c}}{z^{2c+2}} ,
    \qquad \text{convergent iff } D < z ,
    \label{eq:series}
\end{align}
so that, contracted with the vertex,
\begin{align}
    \Qcal^{(2\rho+1)}(z) \;=\; \sum_{c \ge 0} \frac{(-1)^c}{z^{2c+2}}\, \Dcal^{(2\rho+1+2c)} .
    \label{eq:seriesQ}
\end{align}
This is entirely separable --- every term is one of the polynomial objects
$\Dcal^{(j)}$ --- and needs no integral at all. Its cost is that it requires $\Dcal$ at
powers far \emph{above} the moment order: truncating \eqref{eq:seriesQ} at $C$ terms needs
$\Dcal^{(j)}$ up to $j = 2\rho + 1 + 2C$. In the code this truncation is the parameter
\code{expan}, set to $41$ in \code{thc\_26\_10}, raised to $51$ in
\code{small\_Q\_working\_THC}, and cut back to $12$ in the November refactors. The
corresponding routine is \code{build\_small}.

\subsection{The polynomial branch: the exact finite part}

For the complementary regime the same ratio is used the other way. Polynomial division
gives, exactly and for every pair,
\begin{align}
    \frac{D^{2\rho+1}}{D^2+z^2}
    \;=\; \underbrace{\sum_{k=0}^{\rho-1} (-z^2)^k D^{2\rho-1-2k}}_{\text{separable}}
    \;+\; \underbrace{(-z^2)^{\rho}\, \frac{D}{D^2+z^2}}_{\text{not separable}} ,
    \label{eq:division}
\end{align}
which in the auxiliary space reads
\begin{align}
    \Qcal^{(2\rho+1)}(z) \;=\; \sum_{k=0}^{\rho-1} (-z^2)^k\, \Dcal^{(2\rho-1-2k)}
    \;+\; (-z^2)^{\rho}\, \Qcal^{(1)}(z) .
    \label{eq:divisionQ}
\end{align}
The finite sum is \code{build\_large}. It is exact, cheap and needs no expansion parameter,
but it leaves a remainder --- $\Qcal^{(1)}(z)$ --- which is precisely the object that
motivated the whole exercise, now carrying a prefactor $(-z^2)^{\rho}$ that grows with the
moment order.

\subsection{The Laplace branch: the remainder, and an inner quadrature}

The remainder is reached through the Laplace representation
\begin{align}
    \frac{D}{D^2+z^2} \;=\; \int_0^\infty \cos(zp)\, e^{-Dp} \,\mathrm{d}p ,
    \qquad
    e^{-D_{ia} p} = e^{\epsilon_i p}\, e^{-\epsilon_a p} ,
    \label{eq:laplace}
\end{align}
whose exponential \emph{is} separable. Contracting \eqref{eq:laplace} with the vertex gives
\begin{align}
    \Qcal^{(1)}(z) = \int_0^\infty \cos(zp)\, \mat{M}(p) \,\mathrm{d}p ,
    \qquad
    M_{PQ}(p) = \sqrbrckts{\sum_i X_{iP} X_{iQ} e^{\epsilon_i p}}
                \sqrbrckts{\sum_a X_{aP} X_{aQ} e^{-\epsilon_a p}} ,
    \label{eq:Mp}
\end{align}
an elementwise product of two $O(N^3)$ objects. Together, \eqref{eq:divisionQ} and
\eqref{eq:Mp} are \code{build\_large} plus the correction routine that follows it:
\begin{align}
    \Qcal^{(2\rho+1)}(z) \;=\; \sum_{k=0}^{\rho-1} (-z^2)^k \Dcal^{(2\rho-1-2k)}
    \;+\; (-z^2)^{\rho} \sum_{p} w_p \cos(z p)\, \mat{M}(p) .
    \label{eq:largeQ}
\end{align}

The kernel \eqref{eq:Mp} is correct and is kept. The difficulty is the quadrature over $p$.
The integrand oscillates at frequency $z$ while decaying at rate $D$, so a rule with one
scale cannot resolve it when $z \gg D$: the oscillation is fast and the decay slow, and the
cancellation between half-periods is what has to be captured. Every version therefore
carries a nested optimisation --- \code{optimise\_q\_quad}, which rescales an inner
Gauss--Laguerre rule at \emph{every outer point} by minimising the error of a scalar
integral it can evaluate in closed form. Even so, the inner rule was found to be delicate:
the point count was reduced from $20$ to $14$ with the code comment ``\code{20 starts to
introduce error on e-9}'', an unusual failure direction that indicates the optimisation, not
the resolution, was setting the accuracy.

\subsection{The integration-by-parts family, and what it says about the split}

Integrating \eqref{eq:laplace} by parts twice relates the two branches directly. With
$\mathcal{F} = 2\int_0^\infty \cos(zp) e^{-Dp}\mathrm{d}p$,
\begin{align}
    \mathcal{F} \;=\; \frac{2D}{z^2} - \parth{\frac{D}{z}}^2 \mathcal{F} ,
\end{align}
and iterating gives the family
\begin{align}
    \mathcal{F} \;=\; \frac{2}{z}\sum_{l} (-1)^l r^{2l-1} \;+\; (-1)^{l+1} r^{2l}\,\mathcal{F} ,
    \qquad r = \frac{D}{z} .
    \label{eq:ibp}
\end{align}
Each member of \eqref{eq:ibp} is exact; they differ in how much of the answer sits in the
closed-form prefactor and how much in the integral, and therefore in how much any error in
the integral is amplified. Three values make the point:
\begin{center}
\begin{tabular}{llll}
\toprule
$r$ & prefactor on $\mathcal{F}$ & effect on an integration error & regime \\
\midrule
$0.1$ & $10^{-4}$ & suppressed by four decades & $z \gg D$ \\
$1$   & $1$       & neither suppressed nor amplified; polynomial part vanishes & $z \sim D$ \\
$10$  & $10^{4}$  & amplified by four decades & $z \ll D$ \\
\bottomrule
\end{tabular}
\end{center}
So the partially integrated form should be used up to $r = 1$ and the direct integral at
$r \approx 1$, and neither is safe on the other side. That prescription is a statement about
the individual pair $(i,a)$ --- and it is the sentence on which the previous line of work
stalled, because the prescription is per pair while the machinery is not.

\section{Why the partition was forced}

A hypercontracted builder can evaluate a restricted $\Dcal$ only over a product set. For
$S_i \subseteq \{i\}$ and $S_a \subseteq \{a\}$,
\begin{align}
    \Dcal^{(m)}\sqrbrckts{S_i \times S_a}_{PQ}
    = \sum_{p=0}^{m} \binom{m}{p}
      \sqrbrckts{\sum_{i \in S_i} X_{iP} X_{iQ} (-\epsilon_i)^p}
      \sqrbrckts{\sum_{a \in S_a} X_{aP} X_{aQ} \epsilon_a^{\,m-p}} ,
    \label{eq:rectangle}
\end{align}
and likewise for \eqref{eq:Mp}. The two brackets are independent sums, so the only regions
reachable at $O(N^3)$ are \textbf{rectangles} $S_i \times S_a$.

The region the physics asks for is not a rectangle. With the orbital energies sorted
ascending, the set on which the series \eqref{eq:series} converges is
\begin{align}
    R_\theta = \{(i,a) : \epsilon_a - \epsilon_i < \theta \}
             = \{(i,a) : a < a_{\max}(i)\} ,
    \qquad a_{\max}(i) \text{ non-decreasing in } i ,
\end{align}
a \textbf{staircase}: as $\epsilon_i$ rises the admissible $\epsilon_a$ rises with it. A
staircase with $K$ steps is a union of $K$ rectangles, and can be assembled by
inclusion--exclusion --- for two steps,
\begin{align}
    \bm{1}_{R} = \bm{1}_{I \times A} - \bm{1}_{I_1 \times A_1} ,
    \qquad
    \Dcal\sqrbrckts{R} = \Dcal\sqrbrckts{I \times A} - \Dcal\sqrbrckts{I_1 \times A_1} ,
    \label{eq:inclusion}
\end{align}
which is literally the line
\code{(O$_1$+O$_2$)*(V$_1$+V$_2$) - O$_1$*V$_1$} in \code{get\_D\_exp\_list\_zero}. Each
extra step is another rectangle, another pair of collocation contractions, and another term
whose sign must be tracked. Since the boundary $a_{\max}(i)$ moves with $z$, the
decomposition has to be recomputed along the outer quadrature, and pairs migrate between
branches as $z$ sweeps --- which is what the \code{book} array, \code{d\_above\_one} and
\code{d\_lower} exist to track.

Two thresholds control the whole apparatus and neither is derived: \code{d\_thres}
($\nicefrac{1}{3}$, later $0.4$, then $0.36$) sets where the series is trusted, and
\code{integral\_switch} ($\nicefrac{1}{0.7}$) sets where the integral takes over.

\section{The line of attempts}

\begin{center}
\begin{longtable}{@{}p{1.8cm}>{\raggedright\arraybackslash}p{4.5cm}>{\raggedright\arraybackslash}p{8.9cm}@{}}
\toprule
date & file & what changed \\
\midrule
\endhead
2025-09-01 & \code{take1\_thc\_rpa} & first version; a per-pair \code{book} array tracks
    which $(i,a)$ have migrated between branches \\
2025-10-13 & \code{working\_drpa\_}\allowbreak\code{thc\_code} & \code{book} replaced by index slices;
    \code{build\_large\_1/2/3} for one, two or three rectangles \\
2025-10-15/16 & \code{working THC}, \code{working\_THC\_2} & two independent inner
    quadratures \\
2025-10-22 & \code{pre\_combine\_}\allowbreak\code{integrals\_thc} & integrals combined; the Cholesky
    factorisation of $\mat{Z}$ removed \\
2025-10-26 & \code{thc\_26\_10} & the splits precomputed once into an array indexed by
    quadrature point; each $\Qcal$ checked inline against an exact $O(N^4)$ evaluation \\
2025-10-27 & \code{small\_Q\_working\_THC} & only even orders stored, halving $\Qcal$;
    the $\Dcal$ builder vectorised \\
2025-10-29 & \code{THC\_pre\_one\_}\allowbreak\code{integral\_type} & a \code{numba} kernel for the binomial
    contraction \\
2025-11-06/12 & \code{midway\_refactor\_THC},\newline \code{THC\_take2\_pre\_}\allowbreak\code{cleanup} & the binomial
    split by parity to halve the working set; a \code{tail} branch \\
2026-01-03 & \code{RPA\_solids\_thc\_logic} & the periodic plan: sum over the full
    $n_k^2$ rather than per $\mathbf{q}$, since $\mat{Z}$ differs per $\mathbf{q}$ \\
2026-02-02/04 & \code{rpa\_thc\_solids\_wip},\newline \code{backup\_get\_D\_}\allowbreak\code{mask\_start} & slices
    replaced by boolean masks on the orbital indices --- the honest form of
    \eqref{eq:rectangle} \\
2026-03-04 & \code{RPA\_Solids\_THC\_}\allowbreak\code{Working\_Q} & three mask sets, with an exact reference
    computed alongside for validation; 1461 lines \\
\bottomrule
\end{longtable}
\end{center}

The trajectory is a steady fight with the bookkeeping of \eqref{eq:inclusion}, never with
the physics. The final version is explicit about it: its three mask sets are named
\code{series}, \code{low\_int} and \code{high\_int}, and it carries an $O(N^4)$ evaluation
of $\Qcal$ in parallel with the $O(N^3)$ one so that the two can be compared point by point.

\section{Where it broke}

That inline comparison recorded the failure directly. On an eight-carbon alkane, at
quadrature points where a single region is active the check passes:
\begin{center}
\code{point 2.6834535605444727 / One integral / final True 2.47955322265625e-05}
\end{center}
and at points where the partition is active it does not:
\begin{center}
\code{point 1.6810945450652695 / Two integrals / final False 0.0013410896062850952}
\end{center}
a relative error of $2.6\times10^{-3}$ in $\Qcal^{(10)}$ against $2.5\times10^{-5}$ where no
split occurs --- two decades, and at the boundary rather than in the kernel. The failure is
in \eqref{eq:inclusion} and its moving boundary, not in \eqref{eq:series},
\eqref{eq:divisionQ} or \eqref{eq:Mp}.

\section{What replaces it}

Three changes, of which only the second is new physics.

\subsection{The outer quadrature is already gone}

The zeroth moment in \code{momentGW} is no longer computed by an optimised quadrature over
$z$. A certified rational approximation $x^{-1/2} \approx \sum_j w_j (x+s_j)^{-1}$ on a
rigorous spectral enclosure turns \eqref{eq:compact} into a finite pole sum,
\begin{align}
    \Ncal^{(m)}
    = \sum_j w_j \parth{\mat{I} + 4\Qcal^{(1)}(s_j)}^{-1} \Qcal^{(m+1)}(s_j) ,
    \qquad
    \Qcal^{(k)}(s) = \mat{V}\, \frac{\mat{D}^{k}}{\mat{D}^2 + s}\, \mat{V}^T ,
    \label{eq:poles}
\end{align}
over 18 to 19 poles at the production tolerance. The shifts are a fixed, known set. This
happened on the density-fitted side for its own reasons, and it removes
\code{optimise\_main\_quad} and the scale optimisation the hypercontracted code inherited.

\subsection{The inner quadrature becomes a fit}

The oscillatory rule over $p$ is replaced by a least-squares fit on a set of exponents
$\{t_k\}$ chosen once,
\begin{align}
    \frac{D}{D^2+s} \;\approx\; \sum_k c_k(s)\, e^{-t_k D} ,
    \qquad
    e^{-t_k D} = e^{t_k(\epsilon_i - \mu)}\, e^{-t_k(\epsilon_a - \mu)} ,
    \label{eq:expfit}
\end{align}
made against the true set of gaps with relative-error weighting. Structurally this is
\eqref{eq:Mp} with $w_p \cos(z p) \to c_k(s)$; the separable kernel $\mat{M}$ is unchanged.
What changes is everything about how the coefficients are obtained:

\begin{center}
\begin{tabular}{@{}>{\raggedright\arraybackslash}p{3.2cm}>{\raggedright\arraybackslash}p{5.4cm}>{\raggedright\arraybackslash}p{5.6cm}@{}}
\toprule
& previously & now \\
\midrule
coefficient & $w_p \cos(z p)$, from a rule & $c_k(s)$, from a fit \\
integrand & oscillatory, scale set by $z$ & none \\
grid & re-optimised at every outer point & chosen once, shared by every pole \\
varies with moment order & no & no \\
varies with the pole & yes, through the whole rule & only through $c_k$ \\
error control & inferred from the rule & measured on the gaps that are used \\
\bottomrule
\end{tabular}
\end{center}

Because the exponents are shared, the expensive object $\mat{M}(t_k)$ is built once and
reused by every pole and every moment order --- where previously it was rebuilt for each
$(z, p)$ pair. The exponents are chosen by pivoted QR on the exponential basis rather than
laid out geometrically, which reaches the achievable floor with roughly forty rather than
sixty points, and the search returns the smallest grid meeting a requested residual: 21
exponents for lithium hydride, 32 for water, 36 for ozone, 40 for benzene at $10^{-11}$.

\subsection{The partition is not replaced --- it is deleted}

With \eqref{eq:expfit} valid over the whole range of $D$ at once, there is no $r$-dependent
choice to make and no region to build. The powers of the denominator come from the binomial
expansion, exactly:
\begin{align}
    \Qcal^{(m+1)}(s) \;\approx\; \sum_k c_k(s) \sum_{p=0}^{m} \binom{m}{p}\,
        \mat{O}_p(t_k) \Had \mat{U}_{m-p}(t_k) ,
    \label{eq:routeB}
\end{align}
with $\Had$ the elementwise product and
\begin{align}
    \sqrbrckts{\mat{O}_p(t)}_{PQ} = \sum_i X_{iP} X_{iQ} (-\eb_i)^p e^{t\eb_i} ,
    \qquad
    \sqrbrckts{\mat{U}_q(t)}_{PQ} = \sum_a X_{aP} X_{aQ} \eb_a^{\,q} e^{-t\eb_a} ,
\end{align}
referenced to the chemical potential so that both exponentials lie in $[0,1]$ and every
binomial term is a product of like-signed factors. Equation \eqref{eq:routeB} is
\code{get\_D\_exp\_list} --- the kernel the previous versions already had. It is kept
because, measured against an exact evaluation of the same pole sum, it is the only one of
three candidates that survives to high moment order: the polynomial route
\eqref{eq:divisionQ} loses about one and a half decades per order to the $(-s)^k$
prefactors, and fitting $D^{m+1}/(D^2+s)$ directly saturates, a growing function being
unreachable from decaying exponentials without cancellation that double precision will not
carry.

\subsection{What is now certified}

Writing $\delta$ for the fit residual and $\varepsilon = \max_{ia} |\delta(D_{ia})| /
g_s(D_{ia})$, positive semidefiniteness of $\Qcal^{(m+1)}$ gives
\begin{align}
    \big\| \Delta \Qcal^{(m+1)} \big\|_2 \;\le\; \varepsilon \big\| \Qcal^{(m+1)} \big\|_2
    \qquad \text{for every } m .
\end{align}
One measured scalar bounds the relative spectral error at every moment order. The previous
architecture had no analogous statement: its error was the sum of a truncation in
\eqref{eq:series}, a quadrature error in \eqref{eq:Mp}, and a boundary error in
\eqref{eq:inclusion}, of which only the first was controlled by a parameter anyone set.

\section{Correspondence}

\begin{center}
\begin{tabular}{lll}
\toprule
 & previous & now \\
\midrule
outer integral & optimised Clenshaw--Curtis in $z$ & certified poles $s_j$ \\
non-separable factor & partitioned by $r = D/z$ & fitted once, all $D$ \\
$r \ll 1$ & series \eqref{eq:seriesQ}, \code{expan} $= 12$--$51$ terms & --- \\
$r \gg 1$ & division \eqref{eq:divisionQ} + inner quadrature & --- \\
boundary & inclusion--exclusion over rectangles & --- \\
thresholds & \code{d\_thres}, \code{integral\_switch} & --- \\
inner grid & re-optimised per outer point & shared, selected once \\
separable kernel & $\mat{M}(p)$, eq.~\eqref{eq:Mp} & unchanged \\
powers of $D$ & series or division & binomial \eqref{eq:routeB}, exact \\
error control & three uncontrolled sources & one measured residual \\
\bottomrule
\end{tabular}
\end{center}

The kernel and the accumulation \eqref{eq:accum} survive intact. What goes is the machinery
between them: three branches, two thresholds, the rectangle decomposition, the migration
bookkeeping and the nested optimisation --- roughly nine hundred of the fourteen hundred
lines of the last version.

\section{Status}

The scalar layer of \eqref{eq:expfit} --- candidate exponents, pivoted selection, the
per-pole fit, the search for the smallest grid meeting a residual, and the diagnostics ---
is implemented and tested. The density--density and self-energy stages of the
hypercontracted dTDA have been made cubic in operations and quadratic in memory, which is
the requirement the whole construction exists to satisfy. What remains for the dRPA is the
assembly \eqref{eq:routeB} against the collocation matrices and the solver that drives
\eqref{eq:poles} and the moment recursion.

One caution carried forward from the record above. The previous line of work was not wrong
about the kernel, and it was not short of effort; it was defeated by a boundary that had to
be maintained. Any future proposal that reintroduces a per-pair choice --- a threshold, a
mask, a regime --- should be measured at the boundary before it is measured anywhere else,
because that is where the two decades were lost.

\end{document}
